\begin{problem}[33届物理竞赛复赛]{57}
    某秋天清晨，气温为 $4.0^{\circ}$ C，一加水员到实验园区给一内径为2.00 m、高为2.00 m的圆柱形不锈钢蒸馏水罐加水。罐体导热良好。罐外有一内径为4.00 cm的透明圆柱形观察柱，底部与罐相连（连接处很短），顶部与大气相通，如图所示。加完水后，加水员在水面上覆盖一层轻质防蒸发膜（不溶于水，与罐壁无摩擦），并密闭了罐顶的加水口。此时加水员通过观察柱上的刻度看到罐内水高为1.00 m。

    \begin{subproblem}
        从清晨到中午，气温缓慢升至 $24.0^{\circ} \mathrm{C}$ ，问此时观察柱内水位为多少？假设中间无人用水，水的蒸发及罐和观察柱体积随温度的变化可忽略。
    \end{subproblem}
    \begin{subproblem}
        从密闭水罐后至中午，罐内空气对外做的功和吸收的热量分别为多少？求这个过程中罐内空气的热容量。
    \end{subproblem}

    已知罐外气压始终为标准大气压 $p_{0}=1.01\times10^{5}Pa$ ，水在 $4.0^{\circ}C$ 时的密度为 $\rho_{0}=1.00\times10^{3}kg\cdot m^{-3}$ ，水在温度变化过程中的平均体积膨胀系数为 $\kappa=3.03\times10^{-4}K^{-1}$ ，重力加速度大小为 $g=9.80m\cdot s^{-2}$ ，绝对零度为 $-273.15^{\circ}C$ 。
\end{problem}